% ============================================================
% Supplementary File S4: Second-Rater Protocol
% ============================================================

\newpage
\section*{Supplementary File S4: Second-Rater Protocol --- U-T-I-O Case Ratings}
\label{sec:S4}

\noindent\textbf{Purpose:} Mitigate the single-rater limitation on the case-level U-T-I-O ratings
(4~cases $\times$ 4~domains = 16 ordinal ratings on a four-band scale: H~/ M~/ L--M~/ L).
Distinct from the 17-source domain-presence self-audit (audit-check.md), which audits
\textit{domain-presence} coding (Y/blank) on the 17-source corpus.

\noindent\textbf{Scope:} Applies the rubric in Supplementary File~S3 to the four cases
(U.S.\ ONC HIE, NHS Shared Care Records, Estonia eHealth, U.S.\ VA/DoD EHR Modernization).
Each case-domain cell is re-rated by a second rater blind to the original ratings,
then compared and adjudicated.

\bigskip

% ------------------------------------------------------------------
% 1. METHOD
% ------------------------------------------------------------------

\subsection*{1. Method}

\subsubsection*{1.1 Rater Independence}

The second rater receives: (a) the construct-level criteria for U, T, I, and O (the four
rubric tables in Supplementary File~S3); (b) the case-by-case evidence anchors with
citation keys, but \textbf{with the original ratings redacted}; and (c) the bibliography
for source verification. The second rater does \textbf{not} see the original H / M / L--M / L
assignments before producing their own.

\subsubsection*{1.2 Rater Types Reported}

This protocol supports two rater types, reported separately:
\begin{enumerate}
\item \textbf{Human second rater} (preferred; not yet recruited). A credentialed reviewer
  with health-IT or implementation-science expertise re-rates from the blinded rubric.
  Status: pending.
\item \textbf{AI-assisted second rater} (supplementary; executed). Anthropic Claude Opus~4.7,
  accessed 2026-04-29. Given the blinded rubric and access to the cited primary sources via
  web tools. \textbf{Does not substitute} for human dual-rating; mirrors the AI-assisted
  pattern used in the 17-source domain-presence self-audit.
\end{enumerate}

\subsubsection*{1.3 Agreement Statistics}

With N~=~16 ordinal cells on a 4-band scale, three statistics are reported jointly:
\begin{itemize}
\item \textbf{Percent exact agreement} --- fraction of cells where second rater matches
  original exactly (primary; easy to interpret).
\item \textbf{Percent within-one-band agreement} --- fraction where ratings are within one
  band of each other (catches close disagreements that may reflect rubric-band ambiguity).
\item \textbf{Linear-weighted Cohen's~$\kappa$} --- ordinal-aware agreement statistic.
  Weights: $w_{ij} = |i-j|/3$ where $i$, $j$ are band indices on a 1--4 scale
  (L=1, L--M=2, M=3, H=4). Reported with the explicit caveat that \textbf{N~=~16
  yields a high-variance $\kappa$ point estimate}; reported for transparency, not
  as a robust inferential claim.
\end{itemize}

\subsubsection*{1.4 Disagreement Adjudication}

For each disagreement (exact mismatch), the original rater records: which construct(s)
drove the original rating; which construct(s) the second rater appears to have weighted
differently; and whether the disagreement reflects (a)~different evidence interpretation,
(b)~rubric-band ambiguity, or (c)~a defensible rating revision. If category~(c), the
original rating is updated in Supplementary File~S3 with a logged change.

\bigskip

% ------------------------------------------------------------------
% 2. BLINDED RUBRIC
% ------------------------------------------------------------------

\subsection*{2. Blinded Rubric Handed to the Second Rater}

\noindent\textit{Note for second rater: Use the construct definitions in Supplementary
File~S3 and the rating decision rules. For each case below, read the cited sources and
assign one of \{H, M, L--M, L\} per domain. Do not consult the evidence-anchored ratings
section of Supplementary File~S3 before completing your ratings.}

\medskip
\noindent\textbf{Case A --- U.S.\ ONC HIE.}
Primary sources: \textit{onc2025reports}, \textit{holmgren2023policyhie},
\textit{eden2016barriers}, \textit{kruse2016adoption}.
Context: Federated, market-based HIE infrastructure under HITECH Act incentives
and 21st Century Cures Act / TEFCA. Rate U~/~T~/~I~/~O.

\medskip
\noindent\textbf{Case B --- NHS Shared Care Records (UK).}
Primary sources: \textit{nhs2025sharedcare}, \textit{fennelly2020national}.
Context: Centrally coordinated national interoperability programme within a single-payer
health system. Rate U~/~T~/~I~/~O.

\medskip
\noindent\textbf{Case C --- Estonia eHealth.}
Primary sources: \textit{estonia2026ehealth}, \textit{europeancommission2016estoniaehr},
\textit{holmgren2023policyhie}.
Context: Legally mandated national digital health infrastructure on the X-Road exchange
platform, operational since 2008. Rate U~/~T~/~I~/~O.

\medskip
\noindent\textbf{Case D --- U.S.\ VA/DoD Federal EHR Modernization.}
Primary sources: \textit{gao2025vaehr}, \textit{gao2024dodehr}.
Context: Large-scale federal EHR replacement (Oracle Health~/~Cerner; MHS GENESIS)
under FY2018 NDAA mandate. Rate U~/~T~/~I~/~O.

\medskip

\noindent\textbf{Blank rating sheet:}

\renewcommand{\arraystretch}{1.25}
\begin{table}[H]
\footnotesize
\centering
\begin{tabular}{lcccc}
\toprule
\rowcolor{headerblue}
\textcolor{white}{\textbf{Case}} &
\textcolor{white}{\textbf{U}} &
\textcolor{white}{\textbf{T}} &
\textcolor{white}{\textbf{I}} &
\textcolor{white}{\textbf{O}} \\
\midrule
\rowcolor{rowgray}
A --- U.S.\ ONC HIE &  &  &  &  \\
B --- NHS Shared Care Records &  &  &  &  \\
\rowcolor{rowgray}
C --- Estonia eHealth &  &  &  &  \\
D --- VA/DoD EHR &  &  &  &  \\
\bottomrule
\end{tabular}
\end{table}

\bigskip

% ------------------------------------------------------------------
% 3. RESULTS
% ------------------------------------------------------------------

\subsection*{3. Results}

\subsubsection*{3.1 AI-Assisted Second Rater (Anthropic Claude Opus~4.7, 2026-04-30)}

The AI rater was given the rubric criteria (§§U--O of Supplementary File~S3) and the
blinded case definitions in~§2 above. The rater operated in a fresh agent context with no
exposure to the original ratings.
\textbf{Audit-trail note:} A first AI pass was discarded after the rater reported
inadvertent exposure to original ratings in an earlier draft. The ratings below are from
the clean second pass; the discarded pass is not retained.

\medskip
\renewcommand{\arraystretch}{1.3}
\setlength{\tabcolsep}{4pt}

\begin{longtable}{p{2.8cm} p{0.8cm} >{\centering\arraybackslash}p{1.4cm} p{9.6cm}}

\caption{AI-Assisted Second-Rater Ratings and Justifications}
\label{tab:s4-airater}\\

\toprule
\rowcolor{headerblue}
\textcolor{white}{\textbf{Case}} &
\textcolor{white}{\textbf{Dom.}} &
\textcolor{white}{\textbf{Rating}} &
\textcolor{white}{\textbf{Justification ($\leq$2 sentences; constructs cited)}} \\
\midrule
\endfirsthead

\caption[]{AI-Assisted Second-Rater Ratings (continued)}\\
\toprule
\rowcolor{headerblue}
\textcolor{white}{\textbf{Case}} &
\textcolor{white}{\textbf{Dom.}} &
\textcolor{white}{\textbf{Rating}} &
\textcolor{white}{\textbf{Justification ($\leq$2 sentences; constructs cited)}} \\
\midrule
\endhead

\bottomrule
\multicolumn{4}{r}{\footnotesize\textit{Continued on next page\ldots}}\\
\endfoot

\bottomrule
\endlastfoot

\rowcolor{rowgray}
A --- U.S.\ ONC HIE & U & \cellcolor{lowred}\textcolor{lowredtext}{\textbf{L--M}} &
Usability and workflow integration mixed across the federated HIE landscape with
documented clinician burden (\textit{eden2016barriers}, \textit{kruse2016adoption});
Patient engagement via portals exists but adoption is variable and Training adequacy
varies widely across participating organisations. \\

A --- U.S.\ ONC HIE & T & \cellcolor{medamber}\textcolor{medambertext}{\textbf{M}} &
Syntactic interoperability advancing under FHIR/USCDI and the API/exchange layer is
federated (TEFCA QHINs, eHealth Exchange) and operational, but Semantic interoperability
and Data completeness remain voluntary and uneven --- fits ``federated/standards-based
but voluntary~$\rightarrow$~M''. \\

\rowcolor{rowgray}
A --- U.S.\ ONC HIE & I & \cellcolor{medamber}\textcolor{medambertext}{\textbf{M}} &
Strong Interoperability mandate (Cures Act information-blocking) and Incentive alignment
(HITECH/MIPS) plus HIPAA Privacy law alignment are in place, but no unified statutory
mandate to participate in HIE and Information governance varies state-by-state. \\

A --- U.S.\ ONC HIE & O & \cellcolor{medamber}\textcolor{medambertext}{\textbf{M}} &
Leadership commitment at federal ONC level is sustained and Implementation planning
documented across multi-year reports to Congress, but Stakeholder engagement and
Change management at the provider level are uneven. \\

\rowcolor{rowgray}
B --- NHS Shared Care & U & \cellcolor{medamber}\textcolor{medambertext}{\textbf{M}} &
Workflow integration and Trust improving as ShCRs embed into clinician workflows
across ICSs, but Usability scores are mixed across regional implementations and
Training adequacy varies across the 42 ICS footprints. \\

B --- NHS Shared Care & T & \cellcolor{medamber}\textcolor{medambertext}{\textbf{M}} &
Syntactic interoperability uses FHIR UK Core and Semantic interoperability leverages
SNOMED CT, with a federated regional API/exchange layer operational, but Data completeness
and System quality vary across ICS-level instances rather than a single national backbone. \\

\rowcolor{rowgray}
B --- NHS Shared Care & I & \cellcolor{highgreen}\textcolor{highgreentext}{\textbf{H}} &
National digital health strategy is documented and funded (NHS Long Term Plan, ShCR
programme), Privacy law alignment is strong (UK GDPR + Data Protection Act + Caldicott),
and Information governance is mature with national IG frameworks. \\

B --- NHS Shared Care & O & \cellcolor{medamber}\textcolor{medambertext}{\textbf{M}} &
Leadership commitment at NHS England and Implementation planning are documented, but
Stakeholder engagement and Change management vary by ICS and Resource capacity is
acknowledged as constrained (\textit{fennelly2020national}). \\

\rowcolor{rowgray}
C --- Estonia & U & \cellcolor{medamber}\textcolor{medambertext}{\textbf{M}} &
Patient engagement is genuinely high (portal usage well above 50\%) and Trust among
clinicians is generally favourable, but published Usability scores and Training adequacy
detail are limited in the cited sources, preventing a confident H across all five
constructs. \\

C --- Estonia & T & \cellcolor{highgreen}\textcolor{highgreentext}{\textbf{H}} &
National X-Road API/exchange layer is the canonical example of a national exchange
backbone, with enforced Semantic and Syntactic interoperability standards, mandated
Data completeness via the national EHR obligation, and documented high System quality
--- meets the H decision rule cleanly. \\

\rowcolor{rowgray}
C --- Estonia & I & \cellcolor{highgreen}\textcolor{highgreentext}{\textbf{H}} &
Statutory Regulatory mandate (national EHR law since 2008) plus a documented multi-year
national digital health strategy, mature Information governance with patient opt-out,
GDPR-aligned Privacy law, and an enforced Interoperability mandate --- satisfies
``statutory mandate + enforced governance + national strategy~$\rightarrow$~H''. \\

C --- Estonia & O & \cellcolor{medamber}\textcolor{medambertext}{\textbf{M}} &
Leadership commitment and Implementation planning sustained over 15+ years with stable
Vendor governance and Resource capacity, but cited sources provide limited evidence on
Change management and Issue management practices at the operator level. \\

\rowcolor{rowgray}
D --- VA/DoD & U & \cellcolor{lowred}\textcolor{lowredtext}{\textbf{L}} &
GAO findings document a persistent Usability crisis (clinician dissatisfaction at
deployed VA sites), Workflow disruption with productivity loss, and Training adequacy
gaps --- meets the explicit ``Documented persistent user-satisfaction crisis
(e.g., GAO findings)'' trigger for L. \\

D --- VA/DoD & T & \cellcolor{lowred}\textcolor{lowredtext}{\textbf{L--M}} &
Oracle Health/Cerner and MHS GENESIS use HL7/FHIR standards (not ``no common standards''),
but System quality issues (outages, patient-safety incidents) and Data completeness /
Integration burden problems documented in GAO reports indicate substantial
integration debt~$\rightarrow$~L--M. \\

\rowcolor{rowgray}
D --- VA/DoD & I & \cellcolor{medamber}\textcolor{medambertext}{\textbf{M}} &
Statutory direction via FY2018 NDAA and federal Privacy law alignment (HIPAA + federal
records statutes) provide strong policy scaffolding, but Information governance and
Interoperability mandate enforcement between VA, DoD, and community providers remain
partially realised per GAO --- fitting strong-incentives-without-full-mandate M. \\

D --- VA/DoD & O & \cellcolor{lowred}\textcolor{lowredtext}{\textbf{L}} &
GAO documents severe execution failures --- cost overruns, schedule slips, deployment
pauses, Vendor governance disputes with Oracle, and audit findings of inadequate
Change management, Issue management, and Resource capacity --- meets the L decision
rule for ``documented severe execution failures''. \\

\end{longtable}

\subsubsection*{3.2 Comparison with Original Ratings}

\begin{table}[H]
\footnotesize
\centering
\setlength{\tabcolsep}{5pt}
\renewcommand{\arraystretch}{1.2}
\begin{tabular}{p{2.8cm} >{\centering\arraybackslash}p{0.8cm} >{\centering\arraybackslash}p{1.3cm} >{\centering\arraybackslash}p{1.8cm} >{\centering\arraybackslash}p{1.5cm} >{\centering\arraybackslash}p{2.0cm}}
\toprule
\rowcolor{headerblue}
\textcolor{white}{\textbf{Case}} &
\textcolor{white}{\textbf{Dom.}} &
\textcolor{white}{\textbf{Original}} &
\textcolor{white}{\textbf{AI Rater}} &
\textcolor{white}{\textbf{Exact}} &
\textcolor{white}{\textbf{Within-one}} \\
\midrule

\rowcolor{rowgray}
A --- U.S.\ ONC HIE & U & M & L--M & No & Yes \\
A --- U.S.\ ONC HIE & T & H & M & No & Yes \\
\rowcolor{rowgray}
A --- U.S.\ ONC HIE & I & H & M & No & Yes \\
A --- U.S.\ ONC HIE & O & M & M & \textbf{Yes} & Yes \\
\rowcolor{rowgray}
B --- NHS Shared Care & U & M & M & \textbf{Yes} & Yes \\
B --- NHS Shared Care & T & H & M & No & Yes \\
\rowcolor{rowgray}
B --- NHS Shared Care & I & H & H & \textbf{Yes} & Yes \\
B --- NHS Shared Care & O & M & M & \textbf{Yes} & Yes \\
\rowcolor{rowgray}
C --- Estonia & U & H & M & No & Yes \\
C --- Estonia & T & H & H & \textbf{Yes} & Yes \\
\rowcolor{rowgray}
C --- Estonia & I & H & H & \textbf{Yes} & Yes \\
C --- Estonia & O & H & M & No & Yes \\
\rowcolor{rowgray}
D --- VA/DoD & U & M & L & No & \textbf{No} \\
D --- VA/DoD & T & M & L--M & No & Yes \\
\rowcolor{rowgray}
D --- VA/DoD & I & H & M & No & Yes \\
D --- VA/DoD & O & L--M & L & No & Yes \\

\bottomrule
\end{tabular}
\vspace{4pt}\\
{\footnotesize\textit{Note: Original ratings are pre-revision values, as methodologically appropriate for inter-rater reliability computation. Two cells (D-U and D-T) were subsequently revised; see §4.}}
\end{table}

\subsubsection*{3.3 AI-Pass Agreement Tally}

\begin{table}[H]
\footnotesize
\centering
\setlength{\tabcolsep}{8pt}
\renewcommand{\arraystretch}{1.3}
\begin{tabular}{p{7.5cm} p{7.5cm}}
\toprule
\rowcolor{headerblue}
\textcolor{white}{\textbf{Statistic}} &
\textcolor{white}{\textbf{Value}} \\
\midrule
\rowcolor{rowgray}
N & 16 \\
Exact agreement & 6/16 (37.5\%) \\
\rowcolor{rowgray}
Within-one-band agreement & 15/16 (93.8\%) \\
Linear-weighted Cohen's~$\kappa$ & 0.27 (Landis--Koch: ``fair'') \\
\rowcolor{rowgray}
Marginal distribution --- original & L=0, L--M=1, M=6, H=9 \\
Marginal distribution --- AI rater & L=2, L--M=2, M=9, H=3 \\
\rowcolor{rowgray}
Direction of disagreements & All 10: AI more conservative (lower band) \\
Cells $>$1 band apart & 1/16 (D --- VA/DoD U: M vs.\ L) \\
\bottomrule
\end{tabular}
\end{table}

\noindent\textbf{Reproducibility note:}
$\kappa$ computed with linear weights $w_{ij} = |i-j|/3$ over the four-band ordinal
scale (L=1, L--M=2, M=3, H=4). At N=16 the $\kappa$ point estimate is high-variance;
the within-one-band percentage is the more interpretable agreement signal.

\medskip
\noindent\textbf{Interpretation:} The AI rater was systematically more conservative than
the original across all 10 disagreements. The marginal-distribution shift---original modal
H, AI modal M---depresses $\kappa$ even though within-one-band agreement is 94\%.
Two cells (D-U, D-T) are flagged in~§4 as candidates for revision because the rubric's
explicit decision-rule trigger language more closely matches the AI's rating.

\subsubsection*{3.4 Human Second Rater}

Status: pending recruitment. The same blinded rubric (§2) and rating sheet apply.

\bigskip

% ------------------------------------------------------------------
% 4. ADJUDICATION LOG
% ------------------------------------------------------------------

\subsection*{4. Adjudication Log}

\noindent For each of the 10 exact mismatches, the disagreement is classified as
\textbf{(a) different evidence interpretation}, \textbf{(b) rubric-band ambiguity},
or \textbf{(c) defensible rating revision}. Disposition: \textbf{kept} (original
retained); \textbf{revised} (original updated); \textbf{flagged} (unchanged pending
external second-rater confirmation).

\renewcommand{\arraystretch}{1.35}
\setlength{\tabcolsep}{4pt}

\begin{longtable}{p{1.0cm} p{1.3cm} >{\centering\arraybackslash}p{0.4cm} p{1.5cm} p{7.5cm} p{2.5cm}}

\caption{Adjudication Log for All 10 Exact Mismatches}
\label{tab:s4-adjudication}\\

\toprule
\rowcolor{headerblue}
\textcolor{white}{\textbf{Cell}} &
\textcolor{white}{\textbf{Orig~$\rightarrow$~AI}} &
\textcolor{white}{\textbf{$\Delta$}} &
\textcolor{white}{\textbf{Class}} &
\textcolor{white}{\textbf{Reasoning}} &
\textcolor{white}{\textbf{Disposition}} \\
\midrule
\endfirsthead

\caption[]{Adjudication Log (continued)}\\
\toprule
\rowcolor{headerblue}
\textcolor{white}{\textbf{Cell}} &
\textcolor{white}{\textbf{Orig~$\rightarrow$~AI}} &
\textcolor{white}{\textbf{$\Delta$}} &
\textcolor{white}{\textbf{Class}} &
\textcolor{white}{\textbf{Reasoning}} &
\textcolor{white}{\textbf{Disposition}} \\
\midrule
\endhead

\bottomrule
\multicolumn{6}{r}{\footnotesize\textit{Continued on next page\ldots}}\\
\endfoot

\bottomrule
\endlastfoot

\rowcolor{rowgray}
A-U & M~$\rightarrow$~L--M & 1 & (b) Rubric ambiguity &
M/L--M boundary is the fuzziest band on the U domain. Original M cites variable
adoption + workflow concerns. AI L--M cites the same evidence and applies the
``two or more constructs at L--M with usability or training gaps documented in
audits'' decision rule. Both are defensible. &
Kept \\

A-T & H~$\rightarrow$~M & 1 & (a) Evidence interp. &
Original weights Cures Act + FHIR API mandate + TEFCA infrastructure as H.
AI weights TEFCA's federated/voluntary structure + non-bound semantic standards
as fitting ``Federated/standards-based but voluntary~$\rightarrow$~M''. &
Flagged --- borderline H/M; Cures Act enforcement is a defensible H-tipping factor \\

\rowcolor{rowgray}
A-I & H~$\rightarrow$~M & 1 & (a) Evidence interp. &
Original weights Cures Act enforcement + HITECH + HIPAA as collectively H per the
rubric's ``enforcement clause.'' AI treats absence of statutory participation mandate
+ state-by-state IG variation as M. &
Kept --- H is consistent with rubric's enforcement clause \\

B-T & H~$\rightarrow$~M & 1 & (a) Evidence interp. &
Original weights national interoperability standards + ShCR platform as H.
AI cites federated regional implementation across 42 ICSs and variable Data
completeness as M. &
Flagged --- borderline H/M; federated regional structure is a genuine H constraint \\

\rowcolor{rowgray}
C-U & H~$\rightarrow$~M & 1 & (a) Evidence interp. &
Original weights $\sim$99\% patient e-Health Record adoption + clinician advocacy as H.
AI cites limited published Usability/Training data in cited sources as preventing
a five-construct H. &
Kept --- patient engagement and trust evidence supports H per U decision rule \\

C-O & H~$\rightarrow$~M & 1 & (a) Evidence interp. &
Original weights sustained ministry-level governance over 15+ years as H.
AI cites limited evidence on Change management and Issue management at operator
level in cited sources. &
Flagged --- H is supported by execution longevity but construct-level evidence partial \\

\rowcolor{rowgray}
D-U & M~$\rightarrow$~L & 2 & (c) Revised &
Rubric U decision rule states: ``Documented persistent user-satisfaction crisis
(e.g., GAO findings)~$\rightarrow$~L.'' \textit{gao2025vaehr} and \textit{gao2024dodehr}
are GAO findings of persistent dissatisfaction. Original M was more lenient than the
rubric strictly indicates. &
\textbf{Revised: M~$\rightarrow$~L--M} (2026-04-30). Middle ground reflecting
documented crisis plus ongoing remediation; not pure L. \\

D-T & M~$\rightarrow$~L--M & 1 & (c) Revised &
Rubric T decision rule: ``Standards adopted but with substantial integration
debt~$\rightarrow$~L--M.'' GAO documents Oracle/Cerner integration debt, deployment
pauses, system quality issues. AI's L--M tracks the rubric language more directly
than original M. &
\textbf{Revised: M~$\rightarrow$~L--M} (2026-04-30). Aligns with rubric trigger language. \\

\rowcolor{rowgray}
D-I & H~$\rightarrow$~M & 1 & (a) Evidence interp. &
Original weights FY2018 NDAA mandate + congressional appropriations as H.
AI cites partial mandate-enforcement between VA, DoD, and community providers as M. &
Kept --- statutory mandate + sustained appropriations support H \\

D-O & L--M~$\rightarrow$~L & 1 & (b) Rubric ambiguity &
L/L--M boundary on O is fuzzy. Original L--M acknowledges audit findings of
execution gaps. AI L applies ``documented severe execution failures'' trigger ---
GAO reports do document cost overruns, schedule slips, and inadequate management. &
Kept --- but AI's L is defensible; a third rater confirming would tighten the case \\

\end{longtable}

\medskip

\begin{table}[H]
\footnotesize
\centering
\setlength{\tabcolsep}{8pt}
\renewcommand{\arraystretch}{1.3}
\begin{tabular}{lc}
\toprule
\rowcolor{headerblue}
\textcolor{white}{\textbf{Disposition}} & \textcolor{white}{\textbf{Count}} \\
\midrule
\rowcolor{rowgray}
Kept (original retained) & 5 \\
Flagged --- left for external second rater & 3 \\
\rowcolor{rowgray}
Revised (rubric trigger language matched by AI rater) & 2 \\
\midrule
\textbf{Total disagreements} & \textbf{10} \\
\bottomrule
\end{tabular}
\end{table}

\noindent\textbf{Revisions applied (2026-04-30):} Two cells in the VA/DoD case were revised
from M to L--M. The VA/DoD case row was updated from U(M) T(M) I(H) O(L--M) to
U(L--M) T(L--M) I(H) O(L--M) in the rubric (Supplementary File~S3, Case~4) and
propagated to the manuscript results table and conclusion.

\medskip
\noindent\textbf{Important methodological note:} The agreement statistics in~§3.3
(6/16 exact, 15/16 within-one-band, $\kappa~=~0.27$) are computed on the
\textit{pre-revision} ratings. Recomputing post-revision would be circular and is not reported.

\bigskip

\noindent\textit{Created: 2026-04-29 --- Last updated: 2026-04-29}
